\documentclass{article}
\usepackage[utf8]{inputenc}
\usepackage[spanish]{babel}
\usepackage{geometry}
\usepackage{xcolor}
\geometry{a4paper, margin=1in}

\title{Guía de Diseño Visual y Prompts: Meta-Pathfinder}
\author{Departamento de Producto \& UX}
\date{Mayo 2026}

\begin{document}

\maketitle

\section{Paleta de Colores Corporativa}
El diseño de Meta-Pathfinder se basa en una paleta que equilibra la serenidad del aprendizaje con la profundidad de la inteligencia artificial. Se aplica la ley de proporciones 60-30-10 para una armonía visual óptima.

\begin{itemize}
    \item \textbf{Primario (Acento/Enfoque):} \texttt{\#674bb5} (Deep Purple). Representa la intuición y la sabiduría artificial.
    \item \textbf{Secundario (Éxito/Calibración):} \texttt{\#006b5e} (Emerald Teal). Representa la precisión y el estado de flujo.
    \item \textbf{Fondo/Superficie (Base):} \texttt{\#F4FAFF} (Cloud White). Proporciona claridad cognitiva y descanso visual.
    \item \textbf{Textos/Sombras:} \texttt{\#0C1E26} (Midnight Blue). Para legibilidad técnica y contraste.
\end{itemize}

\section{Prompts para Generación de Activos (IA)}

\subsection{Logo del Sistema (DALL-E 3 / Midjourney)}
\textbf{Concepto:} Minimalismo abstracto que combine el cerebro humano con una trayectoria de aprendizaje (camino).

\begin{quote}
    \textit{``Minimalist vector logo for a cognitive AI system named 'Meta-Pathfinder'. The design features a stylized human brain silhouette merging into a fluid, winding path. Flat design, geometric shapes. Colors: Deep Purple (\#674bb5) and Teal (\#006b5e). White background, professional, tech-startup aesthetic, clean lines, high symmetry.''}
\end{quote}

\subsection{Fondo para Inicio de Sesión (Landing Page)}
\textbf{Concepto:} Abstracción de ``Estado de Flow''. Formas orgánicas y etéreas que sugieran enfoque y claridad.

\begin{quote}
    \textit{``Bento-grid inspired abstract background for a web application. Soft organic waves and translucent geometric shards floating in a vast airy space. Gradient transitions from soft light blue (\#F4FAFF) to deep purple (\#674bb5) accents. 8k resolution, cinematic lighting, glassmorphism textures, clean atmosphere, empty center for UI placement.''}
\end{quote}

\section{Directrices de Aplicación}
\begin{enumerate}
    \item \textbf{Iconografía:} Utilizar la librería \texttt{lucide-react}. Mantener un grosor de trazo (\textit{stroke-width}) de 2px para consistencia con la tipografía \textbf{Lexend}.
    \item \textbf{Tipografía:} Toda la comunicación debe realizarse en \textbf{Lexend}, priorizando los pesos \textit{Black} (900) para títulos y \textit{Bold} (700) para métricas clave.
    \item \textbf{Bordes:} Aplicar radios de curvatura grandes (1.5rem a 2.5rem) para suavizar la percepción de la interfaz técnica.
\end{enumerate}

\end{document}
